% =========================================================================
\section{Introduction: The Phenomenon and the Question}\label{p08:sec:intro}

There is a fact about intimate life so ordinary that it is rarely felt as strange, and so strange that it deserves the attention this paper will give it. At the moments when the subject most wants to be understood by the one it loves, in longing, in avowal, in the making of vows, in the wish to be near, it very often does not say plainly what it means. It reaches instead for verse: for the rhythmic, the cited, the borrowed, the formally constrained. It writes a poem, or quotes one; it sets words to a measure they did not need; it speaks, in the most private of registers, in a language that is conspicuously \emph{not} the language of plain report. The lover who could say ``I miss you'' says instead, with Li Shangyin, that the spring silkworm spins its thread until death. The one who could say ``I hope we are together for a long time'' says, with Su Shi, \cjkfont 但願人長久，千里共嬋娟\normalfont.

Why? What does verse do, in intimacy, that plain speech cannot? The question is not rhetorical and its answer is not obvious. Plainly, the lyric is less efficient than report: it says less, per word, about states of affairs. If the function of language in intimacy were the transmission of information about feeling, verse would be a poor instrument, and its persistence across every literate culture would be a puzzle of collective irrationality. That it is instead felt as the \emph{higher} form, that to put one's love into a poem is to dignify it, and to receive a poem is to be loved more fully than to be informed, suggests that something other than information is at stake, and that the difference between poetry and prose is doing work we do not yet understand.

This excursus takes that work as its object. But it proceeds by a method I want to make explicit at the outset, because it differs from the deductive spine of the earlier papers in this series. \textbf{It does not begin from a thesis.} To open with a claim about what the lyric \emph{is} would be to frame the phenomenon before we have seen it, to make the material illustrate a conclusion rather than generate one. Instead the paper passes, in turn, through the theories of desire, of language and music, of rhythm, of the dialectic of desire, of rhetoric, and of the Real, that bear on the question; and it permits a claim to emerge only at the end, as an arrival. The guiding question, \emph{why does the intimate subject address the beloved in verse?}, is held open throughout. Each section deepens it from a different angle; none is allowed to close it prematurely.

\subsection{The Structural Place of This Paper}

A word on where this paper sits. The main arc of the series treats coupling as a relation between subjects: the joint attention and attentional economy through which subjects are formed (Paper V), the epistemic agency at work in the proposal (Paper VI), the generation of value and the gaze of the Other in the external relational field (Paper VII). The present paper appears, at first, to leave that arc: its object is not the relation between two subjects but a feature of \emph{language}. Yet the appearance is misleading. When the intimate subject turns to the lyric, it does not address the concrete other directly; it routes its address through the symbolic order, through the inherited forms, the prior poems, the shared figures, and reaches the beloved only by way of that detour. The pole of the coupling has shifted, from the concrete other to what Lacan calls the \textit{big Other} (\textit{le grand Autre}), the order of the signifier in which the subject's desire is constituted and to which its speech is addressed even when it seems addressed to a person. The lyric is the site where this shift becomes legible. The paper therefore does not leave the category of higher-order coupling; it examines a limiting case of it, in which one pole of the couple is the symbolic order itself.

\subsection{Method and the Hidden Spine}

The risk of a survey that begins from no thesis is that it becomes a catalogue, a review of theories of desire, then of language, then of rhythm, with no through-line. To guard against this, the paper carries a \emph{hidden spine}: not a thesis, but a tension, established in Section~\ref{p08:sec:desire} and surfacing repeatedly thereafter. It is the opposition between two conceptions of desire that the relevant traditions force into contact: desire as \emph{lack} (Hegel, Lacan), in which desire arises from a constitutive want and seeks an object that would fill it; and desire as \emph{production} (Spinoza, Deleuze), in which desire is not the absence of an object but the positive, affirmative activity of a power of existing. The lyric, I will suggest, is the rare object in which these two conceptions are not merely both applicable but \emph{co-present and irreconcilable}: the poem is at once the mourning of an absence and the sheer overflow of language's productive power. The paper does not resolve this tension. As an excursus, it claims the privilege of stopping at the place where the two conceptions meet, in the lyric, and naming that meeting, rather than legislating a winner. This is stated now not as a result but as a map of the terrain to be crossed.

% =========================================================================
\section{The Genealogy of Desire and Affect}\label{p08:sec:desire}

If the lyric has something to do with desire, and the phenomenology of intimate verse, saturated as it is with longing, insists that it does, then we must first ask what desire is, and we must ask it without assuming the answer. The concept has at least four major lineages, and they do not agree. Setting them side by side is not eclecticism; it is the condition for locating, precisely, the tension that the rest of the paper will carry.

\subsection{The Psychoanalytic Line: From Libido to the Object-Cause}

Freud's account begins in an economics. The psychic apparatus is governed by quantities of libidinal energy that seek discharge; the dream-work and the symptom are the transformations this energy undergoes under repression, and Freud names two of its central operations \emph{condensation} (\textit{Verdichtung}) and \emph{displacement} (\textit{Verschiebung}), the compression of several latent contents into one image, and the transfer of intensity from a charged idea to an innocuous neighbour \parencite{freud1900interpretation}. We will meet these two operations again, transformed, as the rhetorical mechanisms of the lyric (Section~\ref{p08:sec:rhetoric}); their appearance here, at the root, is deliberate.

Lacan reworks this economics into a structural theory of the signifier. His decisive move is to distinguish \emph{need}, \emph{demand}, and \emph{desire} \parencite{lacan1977ecrits}. Need is biological and admits of satisfaction; its object can fill it. Demand is need passed through language and addressed to another, and because every demand is also, beneath its literal content, a demand for love, for the unconditional gift of the other's presence, no particular satisfaction ever meets it fully. Desire is the residue: what is left when the satisfiable content of need is subtracted from the unconditional reach of demand. Desire therefore has, strictly, no object that could satisfy it; it has only a cause, what Lacan names the \textit{objet petit a}, the object-cause of desire, which is not a thing one could obtain but the very gap that the failure of demand opens. And because the subject's desire is structured in and by the field of the signifier, the order of the Other, Lacan can write that ``man's desire is the desire of the Other'': both desire \emph{for} the Other's recognition and desire patterned \emph{after} the Other's desire \parencite{lacan1977ecrits}. This is the strong form of \emph{desire as lack}. Its consequence, which Section~\ref{p08:sec:dialectic} will draw out, is that desire's true ``object'' is its own perpetuation: a desire that found its object would cease to be desire.

\subsection{The Hegelian Line: Desire and Recognition}

Lacan's lack has an upstream source in Hegel, read through Kojève. In the \emph{Phenomenology}, self-consciousness is first \emph{desire} (\textit{Begierde}): it relates to the world by negating it, consuming the object to confirm its own independence \parencite{hegel1807phenomenology}. But the consumed object vanishes, and with it the confirmation; the desiring self discovers that it needs an object that is not destroyed by being negated, another self-consciousness, which can \emph{recognize} it. Human desire, on this account, is at bottom the desire to be desired, the desire for recognition, and it sets in motion the dialectic of master and slave. Two features of this account matter for us. First, desire is intrinsically \emph{negative}: it is a relation to what is not, a restlessness that consumes its objects and is never at rest in them; this negativity is the philosophical ancestor of Lacanian lack. Second, desire is intrinsically \emph{intersubjective and infinite}: its real aim is another desire, and since to be recognized by a recognized other is to want the other's recognition to be itself worth having, the structure does not terminate. The Hegelian provenance explains why the Lacanian formula, desire is the desire of the Other, is not a paradox but a near-tautology: it was built into the concept of desire from the moment desire was defined as the desire to be desired.

\subsection{The Political-Economic Line: The Social Production of Need}

Against the apparent timelessness of the psychoanalytic and Hegelian accounts stands a third lineage that historicizes desire. For the classical political economists and for Marx, needs are not a fixed natural endowment but are themselves produced, socially and historically, alongside the means of their satisfaction. The range of a human being's wants, and the forms in which they are felt as wants, are the sediment of a mode of production and its culture; the ``necessary'' needs of one epoch are the luxuries or the unthinkable excesses of another. This line matters here for two reasons. First, it supplies a corrective to any account that would treat the longing expressed in the lyric as a pure, ahistorical constant: the very forms of intimate desire, romantic love as we know it, the cult of the beloved, the lyric ``I'' that addresses a ``you'', have a history, and the classical Chinese material this paper analyses belongs to a specific and elaborate social formation of feeling. Second, and looking ahead to Section~\ref{p08:sec:real}, the political-economic line gives us the category of the \emph{surplus}, value produced beyond what is consumed, accumulated rather than discharged, which will let us pose the lyric's economic peculiarity: its \emph{non-exchangeability}, the way a poem is precisely what cannot be paraphrased into information without remainder. The structural analogy between the unparaphrasable remainder of the poem, the Lacanian \textit{objet a} as the remainder of demand, and surplus as the remainder of exchange will be drawn, with due caution, in Section~\ref{p08:sec:real}.

\subsection{The Productive Line: Spinoza, Deleuze, and Desire Without Lack}

The fourth lineage denies the founding premise of the first two. For Spinoza, the essence of each thing is its \emph{conatus}, its striving to persevere in its being; desire (\textit{cupiditas}) is conatus become conscious of itself, and it is therefore not a lack but a \emph{power}, an affirmative quantity of activity \parencite{spinoza1677ethics}. Deleuze and Guattari radicalize this into a frontal assault on the psychoanalytic picture: desire, they argue, is not founded on lack, does not aim at an absent object, and is not structured like a want awaiting fulfilment; it is \emph{productive}, a machinic process that connects and flows, and the entire apparatus of lack-castration-Oedipus is a secondary capture of this primary productivity \parencite{deleuze1972antioedipus}. To say desire is production is to say there is no constitutive hole at the centre of the subject from which desire issues; there is, rather, a positive process that prior theory has misdescribed as a hole by viewing it through the lens of what it lacks.

\subsection{The Tension Established}

We now have the spine. On one side, \textbf{desire as lack}: Hegelian negativity and Lacanian \textit{manque}, desire as the relation to what is not, whose object is its own continuation and whose paradigm is the gap. On the other, \textbf{desire as production}: Spinozan conatus and Deleuzian flow, desire as affirmative power, whose paradigm is overflow. These are not two theories of the same thing that might be averaged. They are two ontologies of desire, and they will pull the analysis of the lyric in opposite directions: lack will read the poem as the staging of an absence, an elegy for what cannot be had; production will read the poem as the sheer generativity of language, an overflow that needs no missing object to explain it. I do not resolve the tension here, and I will not resolve it at the end. I mark it, and I let it run.

% =========================================================================
\section{Language, Poetry, Music: Difference and Teleology}\label{p08:sec:lpm}

Before we can ask what rhythm does, we need the place of poetry among the symbolic systems with which it is continuous. I propose to locate language, poetry, and music on a single scale, and then to ask after the \emph{telos}, the toward-which, of each.

\subsection{The Scale of Receding Reference}

Consider the three as systems of organized sound that differ in the weight they place on the \emph{signified}. Ordinary language is maximally referential: its sounds are, in the Saussurean account, arbitrary with respect to their meanings \parencite{saussure1916course}, and the whole apparatus is built to carry the signified across, with the material of the signifier, the phonic substance, ideally transparent, a window one looks through rather than at. Music sits at the opposite pole: it has, in any ordinary sense, no signified at all; it does not refer, or refers only by convention and association, and its content, if it has one, is exhausted by its form, the organization of the sonic material itself is the whole of what there is to grasp. Music is, in this sense, \emph{pure signifier}: signifier whose play is not in the service of a signified beyond it.

Poetry sits between. It retains reference, a poem is, unlike a sonata, \emph{about} something, and we can paraphrase it, badly, but it systematically thickens the signifier, foregrounding the phonic and formal substance that ordinary language asks us to see through: rhyme, metre, tonal pattern, parallelism, the shape and sound of the word as word. The lyric is language in which the signifier ceases to be transparent and becomes, partly, opaque, in which we are made to attend to the material of speech and not only to its message. On the scale of receding reference, poetry is the middle term: the signified, dominant in prose, begins its retreat; the signifier, suppressed in prose, begins to rise; and music is the limit toward which this movement points, the vanishing of reference into pure sonic form.

\subsection{The Question of Telos}

Each of the three can be asked what it is \emph{for}, and the answers are revealing. Ordinary language is for communication and representation: it is the instrument of the transmission of contents, and it is well-made to the degree that the signifier disappears into the signified. Music, having no signified, cannot be for the transmission of contents; the long tradition that calls it the language of feeling gestures at its telos, not the representation of an emotion but something closer to its \emph{direct presentation}, the production in the listener of an affective movement that is not \emph{about} anything. Music, one might say, does not describe the movement of feeling; it \emph{is} a movement of feeling, transposed into sound.

And poetry? Its telos is precisely its in-betweenness, and this is the hinge of the whole paper. The lyric is neither pure communication nor pure form. It retains enough reference to be \emph{about} the beloved, the longing, the vow, it has not severed itself from the world as music has, but it borrows from music the thickening of the signifier, the insistence that the material of the words matters in itself. Why would intimate speech want to be \emph{both}? Why retain reference (so that it is still \emph{to her}, still \emph{about} this love) while drawing toward music (so that the words become opaque, material, sung)? The provisional shape of an answer, to be earned, not assumed, is that intimacy requires a language that does two incompatible things at once: it must still \emph{mean}, must still be addressed and referential, for it is the address of one person to another; and it must escape the transparency of mere meaning, must become more than information, because what it would convey, love, the beloved's irreplaceability, the longing that has no satisfiable object, is precisely what cannot be transmitted as a content. The lyric's in-betweenness is not a compromise but a solution to a real problem. Sections~\ref{p08:sec:rhythm} and \ref{p08:sec:real} make the solution precise.

% =========================================================================
\section{Rhythm: Four Perspectives and the Repetition of Difference}\label{p08:sec:rhythm}

If the difference between poetry and prose lies in the thickening of the signifier, then rhythm, the periodic organization of the phonic material in time, is the chief instrument of that thickening, and the crux of the paper. I take it up through four lenses, and the fourth, read through Deleuze, is the theoretical centre of gravity of the whole.

\subsection{Functionalist: Rhythm as Mnemonic, Ritual, Coordination}

The oldest answers are functional. Rhythm aids memory: metrically organized speech is far easier to retain and transmit than prose, and in oral cultures the metrical line is the technology of cultural storage, the law, the genealogy, the epic are versified so that they may survive. Rhythm coordinates bodies: the work-song, the rowing-chant, the march synchronize collective labour and movement. Rhythm marks the sacred: incantation, prayer, and spell are rhythmic because rhythm sets speech apart from the profane flow of report and signals that these words are doing something other than informing. Each of these is real, and each tells us something: rhythm is, from the start, the mark of language doing more, or other, than communicating. But functionalism cannot be the whole story, for it explains rhythm by its uses and leaves untouched the question of why \emph{this} organization of sound should have these powers, and why, in intimacy, where memory and coordination and ritual are not obviously at stake, the rhythmic should still be felt as the register proper to love.

\subsection{Structuralist: The Poetic Function}

Jakobson gives the canonical structuralist answer. Language has several functions, and the \emph{poetic function} is the one that focuses on the message for its own sake, on the palpable, material side of the sign. Its mechanism is precise: the poetic function ``projects the principle of equivalence from the axis of selection into the axis of combination'' \parencite{jakobson1960linguistics}. In ordinary speech, equivalence (similarity, the paradigmatic relations among items that could fill a slot) governs \emph{selection}, while contiguity governs \emph{combination} into a sequence. Poetry imposes equivalence onto the sequence itself: syllables are made equivalent to syllables (metre), stresses to stresses, tones to tones (the Chinese \cjkfont 平仄\normalfont\ patterning), line-ends to line-ends (rhyme), whole grammatical structures to one another (parallelism, \cjkfont 对仗\normalfont). Rhythm, on this account, is the periodic structure that results when equivalence is projected onto combination; it is the form that the foregrounding of the signifier takes. This is exact and indispensable, and it gives us the linguistic mechanism of rhythm. But it is, characteristically, \emph{static}: it describes the structure of the rhythmic text as a synchronic pattern of equivalences, and it does not, on its own, tell us why this structure should move us, or what its relation is to time, to the body, to desire.

\subsection{Psychoanalytic: Rhythm as the Return of an Archaic Pulsation}

Psychoanalysis supplies a genetic depth that structuralism brackets. In Kristeva's account, beneath the \emph{symbolic}, the order of grammar, syntax, predication, the order of the signified, lies the \emph{semiotic} (\textit{le sémiotique}), a domain of pre-linguistic pulsion, rhythm, intonation, and bodily drive organized in what she calls, after Plato, the \textit{chora}: the rhythmic, maternal space of the infant's body before the entry into language proper \parencite{kristeva1984revolution}. Poetic language, on this view, is language in which the semiotic erupts through the symbolic: rhythm, assonance, the music of the line are the return, within signifying speech, of this archaic bodily pulsation, the trace of drive in the order of meaning. Rhythm is thus charged with the body and with the earliest relation, the pre-Oedipal bond; to put one's love into rhythm is, at this depth, to route it through the most archaic stratum of one's relation to the other, the stratum of pulse, heartbeat, the rocked body, the maternal voice. The material opacity of the poetic signifier, what structuralism describes from outside as projected equivalence, is, from inside, the pressure of drive against meaning.

\subsection{Existential and Generative: Rhythm as the Repetition of Difference}\label{p08:sec:repdiff}

Here is the paper's pivot. Each account so far treats rhythm as, in some sense, the \emph{return of the same}: the same metre repeated, the same pattern of equivalences, the same archaic pulse. But this is to misdescribe what rhythm is, and Deleuze's \emph{Difference and Repetition} supplies the correction \parencite{deleuze1968difference}. Deleuze distinguishes two repetitions. There is the bare repetition of the identical, the metronome, the mechanical tick, repetition as the return of the same element, and there is a deeper repetition that is the repetition \emph{of difference}: a repetition in which what returns is never the identical but is each time displaced, varied, charged with the new, such that the ``same'' beat is never the same, because it falls in a different place in the gathering series, against a different ground of expectation, bearing the difference that its position in time confers. Rhythm, properly understood, is repetition in the second sense, not the first. A metre is not a metronome. The line that scans is not the mechanical return of identical feet; it is the production, through periodic return, of difference, of expectation and its fulfilment or frustration, of the syncopation that means only against the ground of the regular, of the variation that is audible only because the pattern lets it be heard. Rhythm is the engine by which repetition generates difference; it is, in Deleuze's terms, the temporal form of difference itself.

This reading does three things at once. First, it rescues rhythm from the static synchrony of structuralism: rhythm is not a pattern of equivalences laid out in space but a \emph{movement in time}, a temporal synthesis in which the past beat is retained, the present beat sounded, the next anticipated, and difference produced in the interval. Rhythm is, in this sense, a small machine of time, a demonstration of how the new is generated through return; and this is why I align it with the \emph{existential}, it is the form in which the lyric stages temporality, becoming, the production of the new from the recurrent. Second, it connects rhythm to the \emph{productive} pole of the spine (Section~\ref{p08:sec:desire}): on the Deleuzian reading, rhythm is not the marking of a lack but the affirmative generation of difference, the sheer productivity of the signifier overflowing into time. The poem, read through rhythm-as-difference, is not an elegy for an absent object; it is a generative process, desire as production rather than lack.

And third, this is the crucial complication, and I will not smooth it over, this productive reading collides head-on with the psychoanalytic account of the preceding subsection and with the entire lack-tradition of Section~\ref{p08:sec:desire}. For Kristeva's rhythm is the return of an archaic pulse, charged with the pathos of a lost bond; it is rhythm as the trace of a separation, the mark of a lack. Deleuze's rhythm is the production of difference, owing nothing to any lost object; it is rhythm as affirmation. The same phenomenon, the beating measure of the loved poem, is read by one tradition as the recurrence of a wound and by the other as the generation of the new. I do not adjudicate. I propose, rather, that rhythm is the phenomenon in which lack and production are \emph{co-present and indiscernible}: each return of the beat is at once the re-opening of the gap (the beat that just passed is gone; the next is not yet; the interval is a small recurrent loss) and the production of difference (the next beat, when it comes, is new, displaced, charged). The metre mourns and generates in the same stroke. This is the first full appearance of the spine's tension in the body of the paper, and rhythm is its privileged site. The thesis-to-come, in Section~\ref{p08:sec:conclusion}, will have to answer to this.

\subsection{Rhythm as the Manifestation and Witnessing of the Subject's Existence}\label{p08:sec:rhythmwitness}

The four perspectives above tell us what rhythm \emph{is} (a projected equivalence, an archaic pulse, a repetition of difference) and what it \emph{does} (foregrounds the signifier, charges it with drive, generates the new in time). But there is a further claim, the most important the paper makes about rhythm, that none of the four states outright and toward which all four point: rhythm is the sensuous \emph{manifestation of the subject's own existence}, and, in the intimate scene, its \emph{witnessing}. I take this up on its own because it is the hinge between rhythm as a linguistic device and rhythm as the thing the lover actually wants when the lover wants verse.

Begin from a fact that the structuralist account, which treats rhythm as a synchronic pattern, must bracket: rhythm is lived only in time, and to live it is to be a subject who retains the beat just gone, sounds the beat now, and leans toward the beat not yet arrived. This threefold hold, retention, presentation, protention, is, in the phenomenological tradition, the very structure of inner time-consciousness, the structure by which there is a subject at all rather than a succession of unconnected instants \parencite{husserl1991time}. To follow a rhythm is therefore not to register an external pattern but to \emph{enact} the temporal synthesis that constitutes one as a subject. The beat is not in the poem alone; it is in the subject's holding-together of the poem in time, and that holding-together is the subject's own existing made audible to itself. Heidegger's thought that the being of the subject is care, a being stretched along between having-been and coming-toward, finds in rhythm a small, repeatable, sensuous figure: the measure is the subject's temporal stretch, given as pulse \parencite{heidegger1962being}. To beat time is to exist out loud.

This is why rhythm, of all the features of language, is the one that cannot be paraphrased and cannot be transmitted as content. A signified can be handed over; the temporal synthesis by which a subject lives a rhythm cannot, for it \emph{is} that subject's existing, which is not a content but an event, and an event that must be performed anew by whoever undergoes it. When the metre moves us, what moves us is not information about the poet's feeling but the participation of our own time-consciousness in a movement, the recruitment of our existing into the beat. Rhythm does not represent the subject; it \emph{presents} it, manifests it, brings it into sensuous appearance, in the only medium in which a subject's existing can appear, namely time. Here the productive reading of the previous subsection receives its deepest justification: rhythm produces difference because it is the production of the subject's own continuation in time, the affirmative event of going-on-existing, beat after beat.

And now the intimate turn, which is the reason this matters for a philosophy of intimacy. In the dyad, my existing is not only manifested by the rhythm I make or undergo; it is offered to be \emph{witnessed}. To give the beloved a poem, or to read one with her, is to set going a rhythm that she lives in her time as I live it in mine, so that two temporal syntheses are made to coincide in the same measure, and each becomes the witness of the other's existing. The shared beat is the smallest unit of being-together: not the exchange of a content but the co-enactment of a duration, the mutual attestation that here is a subject, existing, in time, with me. This is what the lover wants and cannot get from plain speech, which transmits a content and leaves the two subjects' times untouched and apart. The lover wants to be \emph{witnessed in his existing} by the one he loves, and to witness hers; and rhythm, because it is existence given as time, is the form in which one subject's being can be laid open to another's attestation. The poem's measure is a request: \emph{beat with me, and so confirm that I am, and that I am with you.} We will see in Section~\ref{p08:sec:bigother} that this witnessing is routed through the big Other, the inherited measures, the prior poems, so that the attestation two lovers give each other is also underwritten by the whole order of those who beat these measures before; the existential function of rhythm and its coupling with the Other are one structure seen from two sides. And we will see in Section~\ref{p08:sec:real} that because the witnessed existing is a continuation in time and not a completed object, the witnessing, like the desire it serves, has no terminus: there is always a next beat, and so always a next confirmation to be given and received.

% =========================================================================
\section{The Dialectic of Desire}\label{p08:sec:dialectic}

The genealogy of Section~\ref{p08:sec:desire} left the lack-tradition with a formula whose consequences must now be drawn: desire's true object is its own perpetuation. This section develops that dialectic, because it is the engine that will drive the account of poetic failure in Sections~\ref{p08:sec:rhetoric} and \ref{p08:sec:real}.

\subsection{The Object That Is Desire Itself}

If desire is the residue of demand, what remains when the satisfiable content of need is subtracted from the unconditional reach of the demand for love, then desire is, by construction, what no object can satisfy. The \textit{objet a} is not an object that desire seeks and might obtain; it is the object-\emph{cause} of desire, the gap itself given a kind of phantom positivity. From this two things follow. First, every empirical object of desire is a stand-in, a placeholder occupying the place of the cause; when obtained, it disappoints, not by accident but by structure, because it was never the cause but only its momentary embodiment, and desire slides on. Second, the deep point, the only ``object'' adequate to desire is desire itself: what desire ``wants,'' beneath every particular want, is to go on desiring, since to attain a final object would be to extinguish desire, and the subject of desire does not, at the deepest level, will its own extinction. Desire wants to want. This is not a pathology but the structure of desire as such.

\subsection{Desire as the Desire of the Other}

The Hegelian inheritance (Section~\ref{p08:sec:desire}) makes this self-relation intersubjective rather than solipsistic. Desire is the desire \emph{of} the Other in the full ambiguity of the genitive: the subject desires the Other's desire (to be desired, recognized, wanted), and the subject desires \emph{according to} the Other's desire (it learns what is desirable from the field of the Other, takes its objects from that field, wants what the Other is taken to want). Desire is therefore never private property; it is constituted in the symbolic field, addressed to it, and patterned by it. This is the precise sense in which the turn to the lyric is a coupling with the big Other (Section~\ref{p08:sec:bigother}): to put desire into the inherited forms of verse is to acknowledge, in the very form of the utterance, that desire was always routed through the Other, that even my most private longing speaks a borrowed tongue and wants according to a desire not first my own.

\subsection{The Negative and the Infinite: The Hegelian Substratum}

That desire cannot terminate is, in Hegelian terms, a consequence of its negativity. Desire relates to its object by negating it, by consuming, surpassing, refusing to rest in it, and the negation of any determinate object only generates the relation anew with respect to a further object; the movement is intrinsically infinite, a ``bad infinity'' of endless succession unless and until it finds an object that is itself a desire, another self-consciousness, in the labour of mutual recognition. The lyric inherits this structure exactly. The poem does not rest in its object; it negates the sufficiency of any plain naming of the beloved (``I miss you'' is never enough, must be surpassed into figure), and it generates, out of that insufficiency, the endless productivity of figuration. The infinity of love poetry, that there is always another poem, that no poem is the last, that the beloved can always be praised again and otherwise, is the aesthetic form of the infinity of desire. We will see in Section~\ref{p08:sec:real} that this is also why the lyric can never be completed and never needs to be: its incompletion is not a defect but the condition of its continuation.

\subsection{The Causation of Desire and the Mechanism of Reproduction}\label{p08:sec:reproduction}

We must be more exact about \emph{how} desire is caused, and the exactness will hand us the bridge to political economy that Section~\ref{p08:sec:real} and Section~\ref{p08:sec:bigother} require. The Lacanian account is not merely that desire lacks an object; it is that desire is \emph{caused} by an object that is itself a remainder, a leftover of a prior operation. When need passes into demand and demand is met, something is always left over, unmet, because the demand was, beneath its content, a demand for the unconditional, which no finite satisfaction supplies; and this leftover, the \textit{objet a}, is not the absence of satisfaction but a \emph{produced surplus}, the by-product of the very operation of satisfying. Desire is caused by its own remainder. This is the decisive point for what follows: desire is not a static hole but a \emph{cycle}, a process that, in operating, produces the very surplus that causes it to operate again. To desire is to enter a loop in which each satisfaction manufactures the residue that reignites the wanting.

Lacan gives this loop an explicitly economic name. The enjoyment that the subject derives not from attaining the object but from the circling pursuit of it he calls \emph{surplus-enjoyment} (\textit{plus-de-jouir}), and he coins the term on the model of, and in avowed debt to, Marx's surplus-value (\textit{Mehrwert}) \parencite{lacan2007seminarxvii}. The parallel is precise and load-bearing. As capital, in Marx, does not consume its surplus but reinvests it to produce more surplus, so that accumulation is a self-sustaining cycle of production and reproduction \parencite{marx1867capital}, so desire does not discharge its surplus-enjoyment but recirculates it, each turn of the loop reproducing the conditions of the next. Desire, in short, has the form of \emph{reproduction}: it is not a one-time want awaiting a one-time satisfaction but a self-reproducing process that secures its own continuation by producing, at every cycle, the remainder that drives the following cycle. The infinity of desire (Section~\ref{p08:sec:dialectic}) is the infinity of a reproductive loop.

This lets us state with new precision what the lyric does, and prepares the central claim of Section~\ref{p08:sec:real}. The poem is the cultural form that takes up this reproductive loop and sustains it deliberately and beautifully. It does not aim to satisfy the longing (to discharge it would be to end it); it aims to keep the loop turning, to produce, at each reading, the surplus that occasions the next saying. Metaphor, metonymy, and rhythm (Sections~\ref{p08:sec:rhetoric}, \ref{p08:sec:rhythm}) are, in this light, not three devices but three \emph{means} to a single end: the reproduction of generativity, the keeping-open of the cycle by which love can be said again. The poem is a reproduction-machine for desire, and, as the next sections argue, the value it reproduces need not be captured and discharged but can be made to circulate, within the dyad and the community of readers, as a sustaining good.

% =========================================================================
\section{Metaphor and Metonymy: The Two Mechanisms of Failure}\label{p08:sec:rhetoric}

We now bring the dialectic of desire down to the level of the line, to the rhetorical operations by which the poem actually works on language. This is the section the whole architecture has been pointing toward, for it is here that the abstract movement of desire becomes the concrete movement of words. The claim is that the two master tropes of poetry are not ornaments but the two forms that desire's movement takes in language, and that each is, in its own way, a mechanism of \emph{approach without arrival}, a failure that is also a function.

\subsection{The Threefold Alignment}

There is a deep correspondence, noticed by Jakobson and theorized by Lacan, among three pairs \parencite{jakobson1956aphasia, lacan1977ecrits}. Jakobson, studying aphasia, found that language disorders cluster along two axes: a \emph{similarity} disorder (loss of the paradigmatic, the capacity for substitution) and a \emph{contiguity} disorder (loss of the syntagmatic, the capacity for combination); and he aligned these two axes with the two master tropes, \emph{metaphor} (substitution on the basis of similarity) and \emph{metonymy} (connection on the basis of contiguity). Freud's two operations of the dream-work, \emph{condensation} and \emph{displacement}, map onto these: condensation, the overlaying of several contents in one image, is a substitutive, metaphoric operation; displacement, the sliding of intensity along a chain of associations to a contiguous term, is a metonymic one. Lacan completes the alignment by reading the unconscious as structured like a language and locating its two mechanisms in the same tropes: metaphor is the structure of the \emph{symptom} (one signifier substituted for another, the repressed returning in a substitute formation) and is the vehicle by which new meaning, and the subject's relation to the signifier ``I'', is produced; metonymy is the structure of \emph{desire} itself, the endless sliding of signification along the signifying chain, the movement from term to contiguous term that never reaches a final signified. The alignment is therefore:

\smallskip
\begin{center}
\begin{tabular}{lll}
\toprule
\textbf{Jakobson} & \textbf{Freud} & \textbf{Lacan} \\
\midrule
Metaphor (similarity) & Condensation & Symptom; production of meaning \\
Metonymy (contiguity) & Displacement & Desire; sliding of the chain \\
\bottomrule
\end{tabular}
\end{center}
\smallskip

\subsection{Metaphor: Approach by Substitution}

Metaphor names the unnameable by putting another name in its place. The beloved is not described but \emph{substituted for}, she is the bright moon, she is the silkworm's unending thread, she is jade and snow and the moon between the clouds. What cannot be said directly (her irreplaceability, the quality that makes her \emph{her}, which is precisely what has no proper name because it coincides with the \textit{objet a}, the unsymbolizable cause of desire) is approached by the violence of substitution, the carrying-over of a name from elsewhere. Metaphor is the vertical, condensing path: it stacks, it overlays, it compresses the unsayable into a single charged figure. And it fails, necessarily, because the substitute is not the thing; the moon is not she; the figure approaches the \textit{objet a} and does not reach it, for the \textit{objet a} has no name that would be the right one. But the failure is productive: out of it comes meaning, the new sense that the metaphor generates, the surplus of signification that was not there before the carrying-over. Metaphor is the mechanism by which the poem produces meaning out of the impossibility of naming.

\subsection{Metonymy: Approach by Sliding}

Metonymy names by moving to the neighbour. Longing is not stated but enacted as a sliding along a chain of contiguous images, the silkworm, then the thread, then the candle, then the ash, then the tears, each term passing the charge to the next, none of them the object, the movement itself being the meaning. This is the horizontal, displacing path, and it is, in Lacan's reading, the very structure of desire: the signification slides, term to contiguous term, and never arrives at a final signified that would arrest it, because there is no such signified, the chain is open, and desire is exactly this openness, this perpetual referral onward. The metonymic poem does not reach its object; it makes the not-reaching audible as duration, as the drawn-out sliding that is the temporal body of longing. The endlessness of the relative-clause, the catalogue, the parallel series in Chinese verse, the way \cjkfont 相思\normalfont\ (longing) is rendered not as a statement but as an unspooling series, is metonymy as the form of desire's non-arrival.

\subsection{Allusion as a Special Metonymy}

Classical poetry, and Chinese poetry pre-eminently, deploys a third device that I read as a special case of metonymy: \emph{allusion}, the use of the prior text, the \cjkfont 典\normalfont. To allude is to name one's present feeling by placing beside it a contiguous fragment of the inherited symbolic order, an earlier poem, a canonical phrase, a legend, and to let the borrowed term carry the charge. This is metonymy operating not within the poem but between the poem and the whole prior order of poems: the present love is connected, by contiguity in the symbolic field, to every prior love that used these words, and is guaranteed by that lineage. Allusion is thus the precise linguistic mechanism of the coupling with the big Other (Section~\ref{p08:sec:bigother}): it pledges the present private feeling to the inherited order and lets that order stand surety for it. When the lover writes \cjkfont 千里共嬋娟\normalfont, he does not coin the feeling; he borrows Su Shi's, and through Su Shi the millennium of moon-gazing that the phrase metonymically drags in its train. The detour through the Other is, formally, an allusion.

\subsection{Both Tropes Approach and Neither Arrives}

The unifying point: metaphor and metonymy are the two, and, Jakobson suggests, the only two, fundamental movements of poetic language, and both are mechanisms of \emph{approach without arrival}. Metaphor approaches the unnameable cause by substitution and fails to coincide with it; metonymy approaches the object by sliding toward it along the chain and fails to halt at it. The two failures are the two faces of the single structure of desire: the vertical failure to name (metaphor, condensation, symptom) and the horizontal failure to arrive (metonymy, displacement, desire). The poem is built out of these two failures, and, this is the thread to Section~\ref{p08:sec:real}, it is built out of them not despite their being failures but \emph{because} they are: only a language that does not reach its object can keep the object as an object of desire, and only such a language is adequate to a love whose object is, in the strict sense, the \textit{objet a} that no word can name and no chain can reach.

% =========================================================================
\section{The Approach to the Real, and the Non-Exchangeable}\label{p08:sec:real}

We can now state, in the vocabulary that has been assembled, what the lyric does, though still as analysis and not yet as the paper's claim, which is reserved for the conclusion. The lyric is the language that moves \emph{toward the Real}. But this formulation is dangerous, and the danger must be met head-on, for the whole rigour of the account depends on getting it right.

\subsection{The Three Registers, and the Real as the Unsymbolizable}

Lacan distinguishes three registers: the \emph{imaginary} (the order of images, identifications, the ego), the \emph{symbolic} (the order of the signifier, language, law, the big Other), and the \emph{real} (\textit{le réel}), which is not ``reality'' but precisely what resists symbolization absolutely, what falls out of the signifier, the impossible that cannot be said \parencite{lacan1977ecrits, lacan1998seminarxi}. The Real is not behind the symbolic as a deeper truth; it is the symbolic's own internal limit, the point at which signification fails, the hole around which the signifying chain circulates. The \textit{objet a} is an object plucked from this register: the unsymbolizable remainder, the cause of desire that has no signifier. And the gap between the Real and the symbolic, the fact that the signifier never captures the thing, that something always falls out, \emph{is} the place where desire is generated (Section~\ref{p08:sec:dialectic}). Desire is the subjective face of the Real--symbolic difference.

\subsection{The Crucial Caveat: Toward the Real, Not Into It}

It is tempting, and several of the paper's earlier formulations leaned this way, to say that the lyric, by thickening the signifier toward music, \emph{reaches} the Real, that rhythm, drawing language toward pure sonic form, breaks through the symbolic to the unsymbolizable beneath. \textbf{This would be an error, and the most important correction the paper makes.} The Real is, by definition, what cannot be symbolized; and rhythm, metre, tonal pattern, rhyme are not the Real but \emph{highly organized signifiers}, they are structure, the most structured language there is, the very acme of the symbolic, not its outside. The poem does not exit the symbolic into the Real. What it does is subtler and is the precise formulation toward which the whole paper has been working: the lyric, within the symbolic, \emph{stages a directed tension toward the Real}, it organizes the signifier so as to point at, gesture toward, circle around, the limit at which signification fails, without and never crossing that limit. Rhythm is not the Real; rhythm is the symbolic's demonstration of its own edge, the signifier made to vibrate at the frequency of its own limit. The poem is a movement toward the Real conducted entirely in the symbolic, an approach that is constitutively an approach and never an arrival, which is exactly the structure of metaphor and metonymy in Section~\ref{p08:sec:rhetoric}, now seen at the level of the whole. This caveat is not a hedge; it is the saving distinction without which the account collapses into mysticism. The lyric does not touch the Real. It is the most exquisite pointing at the Real of which language is capable.

\subsection{Confluence with Rhythm-as-Difference}

This formulation converges with the Deleuzian reading of rhythm (Section~\ref{p08:sec:repdiff}). The ``movement toward a limit'' is precisely a temporal form, and rhythm-as-repetition-of-difference is its temporal body: each return of the beat re-stages the approach (re-opens the interval, the small gap, the not-yet) and produces the difference (the new beat, displaced, charged), so that the poem's rhythm is the perpetual approaching of an edge that is never reached, the temporal enactment of a directed tension that, by structure, does not terminate. Here the two poles of the spine meet once more and refuse to separate: the approach to the Real is the staging of a lack (the limit is never reached; something always falls out; the poem mourns), and it is at the same time the production of difference (each return generates the new; the poem overflows). The lyric is the phenomenon in which desire-as-lack and desire-as-production are the same movement seen from two sides.

\subsection{The Political-Economic Sidelong Glance: Non-Exchangeability}

One more register completes the picture, and it connects this paper to the value-theory of Paper VII. Ordinary language is \emph{exchangeable}: its informational content can be transposed, paraphrased, translated, restated in other words with no essential loss, because what matters is the signified, which survives the change of signifier, this is what it is for language to be a medium of exchange, a currency of contents. The lyric is \emph{non-exchangeable}. A poem cannot be paraphrased without remainder; ``what it says'' cannot be said otherwise, because in the poem the signifier is not transparent to a detachable signified but is itself part of what is meant, the rhythm, the sound, the specific figure are not the vehicle of the content but constitutive of it. This non-exchangeable remainder, the part of the poem that paraphrase cannot capture, is structurally homologous to two other remainders we have met: the \textit{objet a} as the remainder of demand (what is left when the satisfiable is subtracted, Section~\ref{p08:sec:dialectic}), and, in the political-economic register, \emph{surplus} as the remainder of exchange (value produced beyond what circulation accounts for, Section~\ref{p08:sec:desire}).

I draw this as a \emph{structural analogy} and explicitly not as an identity. The three remainders, poetic, libidinal, economic, are not the same remainder, and to collapse them would be the over-unifying gesture this series has consistently refused. What they share is a \emph{form}: in each, a process of circulation or exchange leaves something that the process cannot absorb, and that unabsorbable surplus is, in each case, the locus of value. The poem's non-exchangeability is what makes it a \emph{gift} rather than a communication: in intimate economy, the lyric functions as the pure form of the gift precisely because it cannot be exchanged, cannot be paraphrased into the currency of information, can only be given. This is why a poem given in love is felt as worth more than the information it carries: it is the formal embodiment of the non-exchangeable, of the remainder that escapes the economy of contents, of love's irreducibility to anything that could be transacted. The directed-tension-toward-the-Real and the non-exchangeable-surplus are the same property described in two vocabularies: what the poem approaches and cannot reach (the Real) is the same as what the poem produces and cannot exchange (the surplus, the gift).

\subsection{The Circle of the Good: A Homology with Generative Justice}\label{p08:sec:circle}

The reproduction-loop of desire (Section~\ref{p08:sec:reproduction}) and the gift-character of the poem (just above) together license a claim that connects this paper, by structural homology, to the theory of \emph{generative justice} developed by Ron Eglash \parencite{eglash2016generative}. Eglash's diagnosis of injustice turns on \emph{alienation in the technical sense}: value, whether the labour-value of classical critique, the ecological value of natural systems, or the expressive value of unalienated creative work, is extracted from the network of relations that generates it and siphoned off into accumulation elsewhere, so that the generating network is depleted rather than nourished by its own productivity. Generative justice, correspondingly, names the condition in which value is \emph{not} extracted but is allowed to \emph{circulate back} into the network that produced it, so that production and the conditions of production sustain one another in a loop: a \emph{bottom-up, self-nourishing circulation of unalienated value}.

The homology with the lyric is exact at the level of form, and I press it only that far. Plain communicative speech is, in this register, the \emph{alienating} form of language: it extracts the signified from the signifier, carries the content away, and discards the material body of the utterance as a spent vehicle; its value is realized by being detached and transmitted elsewhere, leaving the act of saying depleted, used up, nothing left to return to. The lyric is the \emph{non-alienating} form. Because its value is non-exchangeable, because it cannot be detached from its own signifier and carried off as content, the value it generates has nowhere to be extracted \emph{to}; it can only remain in, and circulate within, the relation that produced it. The poem given in love is not consumed in the giving; it is kept, returned to, reread, answered with another poem, and at each turn it regenerates the very intimacy that occasioned it. This is a \emph{circle of the good}: a self-nourishing loop in which the expressive value of the verse circulates within the dyad (and within the wider community of readers who keep the poems alive), reproducing rather than depleting the conditions of its own production. The lyric is to the intimate relation what unalienated generative production is to Eglash's community: the form in which value, refusing extraction, feeds back into and sustains the network that made it.

I hold this, in keeping with the discipline of the series, as a homology of \emph{form} and not an identity of substance. The expressive value circulating in a love poem is not the same kind of value as the ecological or labour value in Eglash's cases, and the intimate dyad is not a commons; to flatten these differences would be exactly the over-unification I have refused. What is shared, and it is not trivial, is the structural contrast between a circulation that extracts value out of its generating network (alienation; communication; exchange) and a circulation that returns value into it (generative justice; the lyric; the gift). The poem is generatively just in the precise structural sense that its value cannot be alienated from the relation it serves.

\subsection{The Lyric Presents Generativity, Not Structure}\label{p08:sec:generativity}

We can now formulate the deepest characterization of the mechanism, the one toward which the means analysed above, metaphor, metonymy, rhythm, have all been pointing, and which the structuralist temptation of this whole field most obscures. It is tempting, having identified the devices, to think that what the poem \emph{presents}, what it offers to be grasped, is a \emph{structure}: a determinate pattern of equivalences, a fixed arrangement of figures, an object whose form the reader is to decode. This is an error of the same family as the error corrected in Section~\ref{p08:sec:real} about the Real. The poem does not present a structure. It presents \emph{generativity itself}, the capacity of language to go on generating, and the specific structure of any given poem is merely the occasion, the means, by which that generativity is set in motion in a reader.

Consider what it is to read a poem, as against decoding a message. A message is exhausted when its content is recovered; once decoded, it is finished, and rereading adds nothing, for the structure has yielded its signified and there is no more to extract. A poem is never exhausted, and rereading is not redundant but generative: each reading, in each reader, and in the same reader at different times, sets the figures and the measure to work again and produces meaning anew, never the identical meaning, always displaced, because the reader's interpretive system, the whole field of associations, memories, and prior readings the reader brings, is itself never twice the same. The poem is not a container of a fixed content but a \emph{generator}: it is built so as to keep producing, in the open and varying field of its readers' interpretive systems, an unterminating semiosis. Its value lies not in any structure it has but in its capacity to \emph{sustain generativity} across the indefinite series of its readings. What the poem presents, in the end, is not even a particular language but \emph{the generativity of language as such}, the inexhaustible capacity of the signifier to mean again and otherwise; the individual poem is a window onto that capacity, an occasion for it, not a fixed thing to be possessed.

This reframes everything the paper has said about the means. Metaphor, metonymy, and rhythm are not the \emph{content} of the poem's achievement; they are the techniques by which generativity is kept open, the devices that prevent the poem from collapsing into a decodable message and so closing the loop. Metaphor keeps open the production of new meaning by refusing the proper name; metonymy keeps open the sliding by refusing the final term; rhythm keeps open the subject's living participation by refusing to be a content at all, presenting instead the subject's own continuation in time (Section~\ref{p08:sec:rhythmwitness}). All three are means to the one end: the maintenance of generativity, the keeping-turning of the reproductive loop (Section~\ref{p08:sec:reproduction}), the sustaining of the circle of the good (Section~\ref{p08:sec:circle}). And this is why the specific structure is not, in the end, the main thing. A poem's particular arrangement of figures matters only insofar as it succeeds in launching and sustaining generativity in the interpretive systems of those who receive it; what is to be valued and explained is not the structure on the page but the continuation of generative life that the structure makes possible. The lyric is irreplaceable in intimate coupling not because it encodes a feeling in a clever structure but because it presents, and keeps presenting, the generativity by which a shared life can go on meaning, and a love go on being said.

% =========================================================================
\section{Coupling with the Big Other}\label{p08:sec:bigother}

We can now make explicit the structural claim deferred since the introduction, and thereby re-attach this excursus to the body of the series.

\subsection{The Approach to the Real Runs Through the Symbolic}

The movement toward the Real (Section~\ref{p08:sec:real}) is conducted entirely within the symbolic; the poem points at the limit using the most organized signifiers there are. But the symbolic order is not the subject's private possession: it is the order of the Other, the inherited field of the signifier, the langue that precedes and exceeds any speaker, the accumulated body of prior poems and canonical forms. Therefore the lyric's approach to the Real is necessarily \emph{routed through the big Other}: the subject can stage its directed tension toward the unsymbolizable only by using the symbolic resources, the forms, the metres, the allusions, that belong to the Other and not to it. To write a love poem is to enlist the whole prior order of poems in the service of a present and private feeling; it is to address one's desire to, and through, the big Other.

\subsection{The Detour: Reaching the Beloved Through the Other}

This is the precise sense in which the lyric is a higher-order coupling whose pole has shifted. When the lover addresses the beloved in verse, the immediate addressee of the utterance is not, structurally, the beloved at all; it is the big Other, the symbolic order in which the forms are kept and from which their authority is borrowed. The beloved is reached only by a \emph{detour} through the Other: the lover speaks to the order of the signifier (chooses the form, invokes the prior poem, submits to the metre), and the beloved receives the address as it returns, refracted, from that order. \cjkfont 千里共嬋娟\normalfont\ reaches her because it has first been addressed to Su Shi, to the moon as a millennial signifier, to the whole tradition of the shared moon; she is reached through that detour and would not be reached, not in this way, not with this weight, by the plain ``I hope we last.'' The poem couples the subject with the big Other, and the concrete other is reached as the destination of a message that has gone the long way round, through the symbolic. This is why the address in verse is felt as more, not less, intimate than plain speech: it does not bypass the detour through the Other (no human address can; we are speaking beings); it makes the detour beautiful, deliberate, and shared.

\subsection{Relation to Papers V--VII}

The series' main arc treated coupling between subjects: the formation of the subject in joint attention (V), the epistemic agency of the proposal (VI), the value and the gaze of the Other in the external field (VII). This paper treats a limiting case of the same category: a coupling in which one pole is not a concrete other but the big Other itself, the symbolic order, and in which the concrete other is reached precisely through that coupling. The excursus thus does not depart from the framework of higher-order coupling; it locates its limit. Where the earlier papers asked how two subjects constitute and are just to one another, this one asks how the medium through which they reach one another, language at its most condensed, works, and finds that the medium itself is a third pole, the Other, through whom the two are joined. That is the structural reason an excursus on the lyric belongs in a series on the philosophy of intimacy: the lyric is where the symbolic infrastructure of intimate coupling becomes visible.

% =========================================================================
\section{Case Study: A Functional Typology of Classical Chinese Love Poetry}\label{p08:sec:case}

The architecture must now answer to the material. I classify a body of classical Chinese love poetry not by period or author but by \emph{literary form and illocutionary function}, by what the poem \emph{does} in the relational field, and I show that the resulting classes map onto the mechanisms developed above. The classification is not decoration; it is a test. If the conceptual architecture is sound, the functional kinds of intimate poetry should fall out as the distribution of the mechanisms across the field of intimate speech-acts. Quotations are kept to the analytic minimum; full texts belong to the anthological appendix and not to the argument.

\subsection{Longing: Metonymic Sliding}\label{p08:sec:case-longing}

The poetry of longing addresses an absent object, and its formal signature is the metonymic chain (Section~\ref{p08:sec:rhetoric}): the feeling is not stated but enacted as a sliding along contiguous images that never arrive. Three sub-variants show the mechanism's range.

\emph{Spatial metonymy}, where the signifier of distance or barrier substitutes for the inaccessible person. Li Shangyin's untitled lines, in which meeting is hard and parting equally hard and the spring wind is powerless among the fading flowers, lodge desire in the impassable interval itself; and in the same poet, the spring silkworm whose thread (\cjkfont 丝\normalfont, homophone of \cjkfont 思\normalfont, ``longing'') is spun only at death, and the candle whose tears (wax) dry only when it is ash, give the metonymic chain in its purest form, feeling slides from creature to thread to ash to tear, no term being the object, the sliding itself being the longing. \emph{Temporal metonymy}, where unbroken duration, sleeplessness, the passage of day into night carries the absence: the \cjkfont 古诗十九首\normalfont\ figure of the belt growing looser by the day renders longing as the body's metonymic record of time. \emph{Object-borne metonymy}, where a token, a red bean, a shared river, the moon, takes on and transfers the charge: Wang Wei's red beans of the south, asked to be gathered because ``this thing most stirs longing,'' transfer the whole of desire onto a substitute signifier; Li Zhiyi's lovers at the head and mouth of one long river, drinking the one water, use the contiguity of a single river to bridge a separation that space cannot close. This class is the densest verification of the metonymy-as-desire thesis: longing can be written without end because its object is always at the next link of the chain, never at the present one.

\subsection{Vow: Address to the Big Other}\label{p08:sec:case-vow}

The poetry of the vow does not describe a feeling; it performs a binding, and its formal signature is the address to the big Other as guarantor (Sections~\ref{p08:sec:rhetoric}, \ref{p08:sec:bigother}). \emph{Calling heaven to witness}: the \cjkfont 上邪\normalfont\ of the Han \emph{yuefu} opens by hailing Heaven (\cjkfont 上邪\normalfont, an outright cry to the big Other) and stakes the vow on a catalogue of cosmic impossibilities, mountains worn flat, rivers run dry, summer snow, the meeting of sky and earth, as the only conditions under which the bond could break; the vow is thereby consigned to the permanence of the symbolic order, dissolved only if nature itself were unmade. \emph{The clasped hand}, where a ritual bodily act fixes the symbolic contract: the \cjkfont 诗经\normalfont's \cjkfont 击鼓\normalfont, with its ``in death or in life, the word once given'' and ``I take your hand, with you to grow old,'' makes the clasped hand the bodily signifier of the covenant, the very source-figure of the betrothal covenant. \emph{The wish for one substance}, where the vow takes the form of a wish to abolish the difference between the two: the \cjkfont 孔雀东南飞\normalfont\ pledge in which one will be rock and one the reed, each guaranteeing the other's constancy by the fixity of a thing; and, most profoundly, Bai Juyi's wish, in the \cjkfont 长恨歌\normalfont, to be in heaven two birds that share a wing and on earth two branches grown into one, a wish that, in the same breath, concedes the limit of all such guarantees (``heaven and earth, long-lasting as they are, will end; this regret runs on without end''). The vow's self-knowledge of its own impossibility, in that last instance, makes it the deepest case in the class. Throughout, the structure is that of Section~\ref{p08:sec:bigother}: the vow is addressed not to the beloved but to Heaven, the cosmos, the symbolic order, which is asked to stand surety for the bond, and the beloved is bound through that detour.

\subsection{Avowal: Demand Becoming Desire}\label{p08:sec:case-avowal}

The poetry of avowal stages the conversion of demand into desire (Section~\ref{p08:sec:dialectic}): it is the speech-act of declaring oneself to an other who may not answer, and its formal signature is the unanswered or self-knowing address. \emph{Seeking and not attaining}, the purest form of desire-as-lack: the \cjkfont 关雎\normalfont\ that opens the \cjkfont 诗经\normalfont, with its ``sought and not got, waking and sleeping he longs, tossing and turning,'' is the primal scene of the dialectic of desire, the not-attaining is not the failure of the longing but its very generator; and the \cjkfont 蒹葭\normalfont\ of the \cjkfont 秦风\normalfont, with the one ``in the water's far reach'' who recedes as one pursues upstream, gives the \textit{objet a} its most exact poetic figure, the object forever on the far bank, approached and never reached. \emph{Naked demand and its dignity}: the \cjkfont 越人歌\normalfont's ``the mountain has trees, the tree has boughs; my heart delights in you, and you do not know'' names the structure of avowal as such, the declaration always faces an other whose answer is not guaranteed; Zhuo Wenjun's \cjkfont 白头吟\normalfont, white as mountain snow, bright as the moon between clouds, makes avowal and severance one act, the demand reasserting its dignity at the point of refusal. \emph{Avowal through a medium}, where the declaration is routed through an art-object: Sima Xiangru's \cjkfont 凤求凰\normalfont\ conducts the avowal through the qin, desire reaching the other by the detour of the instrument and the song. The class verifies Section~\ref{p08:sec:dialectic}: avowal is the site where the satisfiable demand (be mine) is transmuted into the unsatisfiable desire (the longing that survives, and indeed feeds on, the uncertainty of the answer).

\subsection{The Wish for Shared Being: The Poetic Suturing of the Real's Gap}\label{p08:sec:case-wish}

The poetry of the wish for shared being is the capstone, and it corresponds to the attempt to suture the Real--symbolic gap (Section~\ref{p08:sec:real}). \emph{Divided yet joined}, where, separation being conceded as irremovable, a shared signifier is offered as the suture: Su Shi's \cjkfont 但願人長久，千里共嬋娟\normalfont, the master specimen, does not deny the thousand miles (the symbolic cannot close that real distance) but offers the shared moon (\cjkfont 嬋娟\normalfont) as the signifier both gaze upon, a suturing that is poetic and imaginary, the gap unclosed, the approach never an arrival; this is the perfect figure of ``toward the Real, not into it'' (Section~\ref{p08:sec:real}), and the poem is candid about it, naming in the same breath that the moon waxes and wanes and ``this was ever hard to keep whole.'' Zhang Jiuling's ``the bright moon is born over the sea; though at the horizon's edge, we share this moment'' gives the same suture in another image, the shared instant bridging the separating distance. \emph{The wish to abolish the boundary}, the most radical suturing of the difference between two beings: Guan Daosheng's \cjkfont 我侬词\normalfont, in which the two are clay kneaded together, ``in my clay there is you, in your clay there is me'', and would share one quilt in life and one coffin in death, wishes to dissolve the very gap between subjects, and yet preserves its own self-knowledge as figure, as clay, as a making rather than a fact. The capstone status of Su Shi's wish is not only thematic: it does not conceal the gap (the miles, the waxing and waning, ``ever hard to keep whole''); it sutures across the gap with a shared signifier while knowing the suture is a figure, which is exactly the structure of the whole paper's account, the directed tension toward a limit that is approached, named, and never crossed.

\subsection{The Reflexive Significance of the Classification}\label{p08:sec:case-reflexive}

The four classes are not an external grid imposed on the poems; they are the distribution of the four mechanisms across the field of intimate speech. Longing is metonymy (Section~\ref{p08:sec:rhetoric}); the vow is address to the big Other (Sections~\ref{p08:sec:bigother}); avowal is the conversion of demand into desire (Section~\ref{p08:sec:dialectic}); the wish for shared being is the poetic suturing of the Real's gap (Section~\ref{p08:sec:real}). That the functional kinds of love poetry should sort themselves exactly onto the mechanisms is the case study's central finding: it shows that the conceptual architecture of the body is not an arbitrary assemblage but tracks a single underlying structure, the structure of desire and its relation to the signifier, of which the genres of intimate verse are the language-level projection. A reflexive observation sharpens this. The four classes lie along a gradient of the object's attainability, and hence of the distance of the \textit{objet a}: avowal (object present, may be sought) $\rightarrow$ longing (object absent) $\rightarrow$ vow (object present, pledged toward an uncertain future) $\rightarrow$ wish for shared being (object permanently divided, yet wished co-present). This gradient traces the recession of the object-cause from near to forever-divided, and it is isomorphic with the approach-to-the-Real of Section~\ref{p08:sec:real}: the genres of love poetry are, read in series, a map of desire's relation to its impossible object.

\subsection{A Note on Boundary Cases}\label{p08:sec:case-boundary}

A typology earns its keep by handling hard cases, and three test the present one. \emph{Elegy for the dead beloved} (Su Shi's \cjkfont 江城子\normalfont\ for his late wife, ``ten years, the living and the dead, each boundless''; Yuan Zhen's ``having once seen the sea, no water is worth the name''): here the object has passed wholly into the Real, it is not absent but unsymbolizable in the strongest sense, death being the Real's absolute figure, and longing reaches its limit, where metonymy slides toward an object that can never, even in principle, occupy the next link. Elegy is thus not a fifth class but the limit of longing, the point at which the metonymic chain confronts an object withdrawn past all recovery. \emph{The poetry of grievance and the abandoned wife} (the \cjkfont 氓\normalfont\ of the \cjkfont 诗经\normalfont; Ban Jieyu's fan laid by in autumn) is the negative of the dialectic of desire: it is what avowal becomes after the answer has come and gone, desire's photographic negative, structured by the same mechanisms in reverse. These boundary cases do not break the typology; they locate its edges, and in doing so confirm that the classification is tracking the mechanisms and not merely sorting themes, for it is precisely at the mechanism-level that the hard cases find their place.

% =========================================================================
% [Inserted section: drafts the author's discussion of the poem as a generative
%  grammar and the symmetry-breaking poetics of truth. Author to review the
%  English rendering before adoption.  <!-- ai-draft --> ]
\section{The Poem as a Generative Grammar, and the Poetics of Truth}\label{p08:sec:grammar}

The case study confirms that the genres of intimate verse are the language-level projection of the mechanisms. We may now press the account of generativity (Section~\ref{p08:sec:generativity}) to its sharpest form, and, in doing so, discover that the method by which this paper has proceeded is not external to its object but is the very thing the object teaches.

\subsection{From Generator to Generative Grammar}\label{p08:sec:grammar-grammar}

Section~\ref{p08:sec:generativity} held that the poem presents not a structure but a generativity, that it is a generator and not a container. We can now say what kind of generator it is. A poem is not a sentence; it is a \emph{grammar}. A sentence is finite and exhaustible: once parsed, it has yielded its content and there is no more. A grammar is finite in its rules and \emph{infinite in its yield}, a small set of rules generating an unbounded set of well-formed sentences, none of which was ``contained'' in the rules and each of which the rules nonetheless license. This is the precise sense of the intuition that the poem ``is a language system'' (Section~\ref{p08:sec:lpm}): the lyric does not deposit in the interpreter a meaning (a sentence) but installs in her a \emph{generative grammar} (a productive capacity), so that what she receives is not a signified but a rule for generating signifieds without end.

This is why rereading is not redundant but generative (Section~\ref{p08:sec:generativity}). To reread a message is to re-parse one finite string and recover the one content again; nothing is produced. To reread a poem is to run, once more, the grammar it installed, against the now-altered ground of the reader's interpretive system, and so to \emph{derive a new sentence}, licensed by the same rules and yet never the one derived before. The poem reads differently each time not because the reader misremembers a string but because she is not re-reading a string at all; she is generating, with an internalized grammar, sentences the poem made possible but did not pre-contain. In the formal-language vocabulary of Paper~IV, the lyric is an internalized generative grammar of growth rate $\lambda > 1$: each derivation, each act of interpretation, produces new and non-repeating yet well-formed output, and the language it generates is, by construction, inexhaustible. What the poem transfers is the \emph{generative rule}, not the generated product, and that is the formal content of its irreducibility to any paraphrase.

And the derivations are not flat. Borrowing, as a structural homology and not an identity, the geometric-phase vocabulary of Paper~VII: each derivation returns not to the same place but carries a \emph{phase increment}, a small sublimation, so that the series of interpretations does not merely repeat but \emph{accumulates}, winding upward. The poem's life in a reader, and across its readers, is the accumulation of this interpretive phase, never a return of the identical, always a charged displacement, which is, at the level of meaning, the very repetition-of-difference that Section~\ref{p08:sec:repdiff} found at the level of the beat.

\subsection{The Generativity That Generates Generativity}\label{p08:sec:grammar-self}

This lets us state, in its purest form, the orientation of production toward production itself that Section~\ref{p08:sec:reproduction} found in the reproductive loop of desire. Ordinary generativity is oriented to perpetuating itself, but it perpetuates by producing some \emph{product}, value, offspring, a work, through which the generating goes on. The poem is more radical, and is perhaps the limit case: what it generates is generativity \emph{itself}. Its product is not a meaning but the capacity to go on generating meaning; it is a system whose output is its own continued power to output. The lyric is generativity become \emph{self-referential}, production whose product is production, the loop closed not upon a thing but upon its own turning.

Here, too, lies the answer to a question the paper has carried without naming: why the lyric is felt as \emph{beautiful}. In other forms, generativity is encumbered, it produces, alongside itself, the bread, the rose, the child, and is seen only through that material product. In the poem, generativity sheds almost all such encumbrance and reproduces itself with the least possible material adhesion: a few borrowed words, a measure, and the thing renews. The lyric is generativity regarding itself in a mirror, the near-pure self-reproduction of the power to mean. Beauty, in this register, is the name for what is witnessed when generativity is seen perpetuating itself in nearly unalloyed form, with the matter through which it works reduced almost to vanishing. That is what one meets in a great line: not a thing said, but the capacity to say, renewing itself before one's eyes.

\subsection{Why a Generator Cannot Be Settled}\label{p08:sec:grammar-settle}

This furnishes, at last, the structural reason for the non-exchangeability of Section~\ref{p08:sec:real}, and connects it to the series' recurring concern with the failure of settlement, of \emph{accounting} (the ledger of Paper~IV). One can settle a product; one cannot settle a generator. Any single output of the poem, this interpretation, the meaning taken on this reading, can be written down, exchanged, balanced in the ledger of contents. But the generator is not located in any one of its outputs; it is the rule that exceeds every output it licenses. To settle a poem is therefore a category error: it is to mistake the generator for one of its products, to take the sentence for the grammar, and so to declare closed an account that is, by its nature, never closed. The lyric's resistance to settlement is not a contingent difficulty but a structural refusal: it cannot be reduced to any of its derivations, and so can never be paid out in full.

The same distinction sharpens the difference between true and counterfeit verse that the paper has needed and not yet stated. The pseudo-lyric is the \emph{idling} generator: formally it too refuses settlement, withholds any decodable content, and strikes the pose of inexhaustibility, but it has installed no working grammar in the interpreter, generates no new sentences, and merely declares, in place, ``I am a generator,'' without the rule that would make the declaration true. In the present notation, true poetry installs a grammar of $\lambda > 1$ with non-zero phase, each derivation new and accumulating; the pseudo-lyric installs a dead grammar of $\lambda = 0$ with zero phase, from which nothing is derived and nothing winds upward, an openness that is not generativity but only the empty gesture of refusing to close. The test of a poem is not whether it withholds a meaning but whether it confers a generativity, whether, in the reader, the rule actually turns.

\subsection{The Poetics of Truth: Interpretation as Symmetry Breaking}\label{p08:sec:grammar-truth}

We arrive at the convergence that justifies, in retrospect, the method announced in the Introduction. Before its reader, a poem's field of meaning is highly \emph{symmetric}: polysemous, suspended, undecided, with many interpretations coexisting as latent and equally licensed, none yet actual. To interpret is to \emph{break} this symmetry, to fall, in this reader and this moment and this situation, onto one determinate reading, and so to let meaning crystallize out of the suspended field. But the responsible interpretation does not mistake its crystallization for the poem. It breaks the symmetry without claiming to exhaust it; it knows that another reader, another hour, will break it elsewhere, and that the poem's meaning is not any single crystallization but the \emph{accumulated phase of the whole series} of symmetry-breaking readings, each partial, each actual, none final. A poem is the standing possibility of an unending sequence of such breaks, and its meaning lives in the sequence, not in any term of it.

This is, exactly, the epistemic posture the series has practiced and that this paper has practiced most openly. The Introduction refused to begin from a thesis and carried instead a tension held open; the conclusions of the series are polyphonic, plural, deliberately unreconciled. We may now see that this is not a stylistic preference but a \emph{poetics of truth}. Before a relational truth, one that no single framework can exhaust because the matter is constituted by relations exceeding any one framework's reach, three postures are available, and only one is responsible. The first is to suspend judgement, to refuse all frameworks in the name of the truth's inexhaustibility, but this is the idling generator, the $\lambda = 0$ openness that produces nothing, the zero-phase pseudo-poem of thought. The second is to declare one framework the whole truth, to settle the account, but this is dogmatism, the category error of taking a single derivation for the grammar. The third is the poetic posture: at this node, in this situation, to break the symmetry onto a single framework, the geometric, the psychoanalytic, the political-economic, the Daoist, and to push its interpretation to the hilt, to say it hard and to the end, while knowing and \emph{marking} that this is one break, one partial crystallization, and then to pass the word to the next. Understanding does not arrive at a resting term; it \emph{spirals upward} through the series of breaks, accumulating phase, as the reading of a poem does.

To hold the conclusions of an inquiry polyphonically, then, each voice partial, each daring to assert, none settling the account, each inviting the next, is to read truth as one reads a poem. The method of this series and the structure of the lyric are not analogous by coincidence; they are one structure, seen once in the order of knowing and once in the order of language. The lyric does not reach the Real, and the inquiry does not reach the whole of its truth; both keep reaching, by a disciplined succession of approaches that are, each one, an approach and never an arrival. That a paper on why we speak love in verse should end by finding that we must also \emph{think} the relational in verse, in the poem's own mode of the unending, partial, accumulating approach, is not a flourish: it is the point at which the object of the paper and its method are revealed to have been, all along, one thing.

% =========================================================================
\section{Conclusion: The Claim Emerges}\label{p08:sec:conclusion}

Only now, having crossed the terrain, may the claim be stated, as an arrival, and in the terms the crossing has earned. The guiding question was: why does the intimate subject address the beloved in verse rather than in plain speech? The answer assembled is this. Plain speech is exchangeable: it transmits a content, and a content about love, about the beloved's irreplaceability, about a longing whose object is the unsymbolizable cause of desire, is precisely what cannot be transmitted, because it is not a content but a relation to a limit. The lyric is the language that, by thickening the signifier toward music through rhythm, stages within the symbolic a directed tension toward the Real: it approaches, by metaphor and by metonymy, the unnameable cause and the unreachable object, and it does so through the inherited forms of the big Other, reaching the beloved by the long and beautiful detour through the symbolic order. The poem is the non-exchangeable remainder, the gift, the surplus that the economy of communication cannot absorb; and that is why it is felt, in love, as worth more than what it says.

And it fails. By the structure of desire (Section~\ref{p08:sec:dialectic}), the approach to the Real is constitutively an approach and never an arrival: metaphor does not coincide with the unnameable, metonymy does not halt at the object, rhythm vibrates at the symbolic's edge without crossing it, the suture of \cjkfont 千里共嬋娟\normalfont\ does not close the thousand miles. But, and here is the claim in its proper form, the failure is the function. Because desire's true object is desire itself, a language that reached its object would extinguish the love it serves; only a language that approaches without arriving keeps the beloved as the object of an unending desire, and so allows love to be spoken without end. The infinity of love poetry, that there is always another poem, is the aesthetic form of the inexhaustibility of a desire that, by structure, cannot be filled. The lyric is irreplaceable in intimate coupling not despite failing to say what it means but because its way of failing is the only way of keeping open the love it speaks.

What the lyric finally presents, then, is not a structure but a generativity (Section~\ref{p08:sec:generativity}). The devices, metaphor, metonymy, rhythm, are means, and the end they serve is the keeping-open of a loop: the reproduction of desire's surplus (Section~\ref{p08:sec:reproduction}), the witnessing of the subject's continued existing in time (Section~\ref{p08:sec:rhythmwitness}), and, in the relation, a circle of the good in which the non-exchangeable value of the verse circulates back into the intimacy that produced it rather than being extracted and spent (Section~\ref{p08:sec:circle}). This is the structural sense in which the lyric is generatively just: its value cannot be alienated from the bond it serves, and so the saying of love nourishes rather than depletes the conditions of its own continuation. A poem is not a message that, once decoded, is done; it is a generator that, reread in the ever-shifting interpretive system of the one who receives it, produces the bond anew at each turn. The particular structure on the page matters only as the occasion of this generativity; what is to be valued is the continuation it makes possible, the going-on of a shared life that keeps meaning, and keeps being said.

I have promised not to dissolve the spine's tension, and I will keep the promise, for the place where it refuses to dissolve is the true location of the claim. Read through \emph{lack}, the lyric is an elegy: it mourns the object that no word can name and no chain can reach, it is structured around a constitutive absence, and its beauty is the beauty of a longing that knows it will not be filled. Read through \emph{production}, the same lyric is an overflow: it is the affirmative generativity of the signifier, rhythm producing difference, language generating the new without need of any missing object to explain its abundance, desire as the sheer power of saying-again-and-otherwise. The poem is both at once, and it is the rare object in which the two ontologies of desire are not merely both true but \emph{indiscernible}: each beat of the metre is at once the re-opening of the gap and the production of difference; each figure is at once the marking of what cannot be named and the generation of new meaning; the wish for shared being is at once the mourning of the thousand miles and the joyous production of the shared moon. I do not choose between lack and production, because the lyric does not. It is the language in which the two faces of desire are shown to be one movement, the movement toward the Real, conducted in the symbolic, destined to fail, and beautiful, and unending, for exactly that reason.

\medskip
{\color{warnred}\itshape Why, then, do we not simply say what we mean to the one we love? Because what we mean has no plain name and lies at the far bank of every word; and because to say it once and be done would be to have finished loving. So we say it in verse instead, borrowing the old forms, beating the old measure, reaching her through the long detour of everything that was ever sung, and we say it again, and otherwise, and without end. The poem does not reach her. It keeps reaching. That is what it is for, and that is what love, in language, is.}
